% Options for packages loaded elsewhere
\PassOptionsToPackage{unicode}{hyperref}
\PassOptionsToPackage{hyphens}{url}
\documentclass[
]{article}
\usepackage{xcolor}
\usepackage{amsmath,amssymb}
\setcounter{secnumdepth}{-\maxdimen} % remove section numbering
\usepackage{iftex}
\ifPDFTeX
  \usepackage[T1]{fontenc}
  \usepackage[utf8]{inputenc}
  \usepackage{textcomp} % provide euro and other symbols
\else % if luatex or xetex
  \usepackage{unicode-math} % this also loads fontspec
  \defaultfontfeatures{Scale=MatchLowercase}
  \defaultfontfeatures[\rmfamily]{Ligatures=TeX,Scale=1}
\fi
\usepackage{lmodern}
\ifPDFTeX\else
  % xetex/luatex font selection
\fi
% Use upquote if available, for straight quotes in verbatim environments
\IfFileExists{upquote.sty}{\usepackage{upquote}}{}
\IfFileExists{microtype.sty}{% use microtype if available
  \usepackage[]{microtype}
  \UseMicrotypeSet[protrusion]{basicmath} % disable protrusion for tt fonts
}{}
\makeatletter
\@ifundefined{KOMAClassName}{% if non-KOMA class
  \IfFileExists{parskip.sty}{%
    \usepackage{parskip}
  }{% else
    \setlength{\parindent}{0pt}
    \setlength{\parskip}{6pt plus 2pt minus 1pt}}
}{% if KOMA class
  \KOMAoptions{parskip=half}}
\makeatother
\usepackage{longtable,booktabs,array}
\newcounter{none} % for unnumbered tables
\usepackage{calc} % for calculating minipage widths
% Correct order of tables after \paragraph or \subparagraph
\usepackage{etoolbox}
\makeatletter
\patchcmd\longtable{\par}{\if@noskipsec\mbox{}\fi\par}{}{}
\makeatother
% Allow footnotes in longtable head/foot
\IfFileExists{footnotehyper.sty}{\usepackage{footnotehyper}}{\usepackage{footnote}}
\makesavenoteenv{longtable}
\setlength{\emergencystretch}{3em} % prevent overfull lines
\providecommand{\tightlist}{%
  \setlength{\itemsep}{0pt}\setlength{\parskip}{0pt}}
\usepackage{bookmark}
\IfFileExists{xurl.sty}{\usepackage{xurl}}{} % add URL line breaks if available
\urlstyle{same}
\hypersetup{
  hidelinks,
  pdfcreator={LaTeX via pandoc}}

\author{}
\date{}

\begin{document}

\section{The Retreat of the False Accusation Plot in Law \& Order: SVU,
1999-2026}\label{the-retreat-of-the-false-accusation-plot-in-law-order-svu-1999-2026}

\textbf{Tyler Satchel Orden} Independent researcher tylerorden@gmail.com

\emph{Manuscript prepared for submission to Television \& New Media}

\begin{center}\rule{0.5\linewidth}{0.5pt}\end{center}

\subsection{Abstract}\label{abstract}

Law \& Order: Special Victims Unit spent two decades as network
television's loudest advocate for believing victims of sexual violence,
while routinely airing plots in which an accusation lands on an innocent
person. Drawing on a census of the show's full run (1999-2026), coded at
the level of the falsely accused person (521 person-episode records
across 567 audited episodes), I compare three eras: pre-\#MeToo (Seasons
1-18), the \#MeToo years (Seasons 19-21), and post-2020 (Seasons 22-27).
The share of episodes containing a false suspect fell from 65.3\% to
50.0\% to 41.8\% (chi-square(2, N = 567) = 20.89, exact p = .000029),
and falsely accused persons per episode fell from 1.06 to 0.57 to 0.56.
The per-season series complicates the era story: Seasons 19 and 20 sit
at or above the preceding plateau, and the break arrives in Season 21,
Warren Leight's first season back as showrunner and the season the
pandemic cut short, so movement, showrunner, and production effects
cannot be separated here. The surviving plots grew harsher during the
\#MeToo-era seasons specifically: severe outcomes rose from 60.0\% of
coded persons to 82.1\% in Seasons 19-21, then settled at 67.3\% after
2020, within the noise band of baseline (p = .018, uncorrected and read
as suggestive, an era-bounded spike rather than a durable shift).
Accusations originating in the victim's own identification roughly
doubled, from 15.5\% to 30.8\%. Depicted police coercion of innocent
suspects traces a U-shape, bottoming at 27.7\% in Seasons 16-18 and
rebounding to 56.7\% and 60.0\% after 2020, though the pooled era test
is nonsignificant (p = .071). Whatever drove the retreat, both the
backlash account and the reform account of post-2017 crime television
fail on the genre's flagship.

\textbf{Keywords:} Law \& Order: Special Victims Unit; \#MeToo; false
accusation; crime drama; content analysis; copaganda; cultivation
theory; longitudinal analysis

\begin{center}\rule{0.5\linewidth}{0.5pt}\end{center}

\subsection{1. Introduction}\label{introduction}

In the last week of September 2017, Law \& Order: Special Victims Unit
began its nineteenth season. Eight days later the New York Times
published its first Harvey Weinstein story, and within a month ``believe
survivors'' had moved from activist shorthand to a mainstream editorial
position. If any fictional institution embodied that position already,
it was SVU's Manhattan Special Victims squad. The show had spent
eighteen seasons dramatizing the claim that victims of sexual violence
deserve to be heard, and its star, Mariska Hargitay, had built a public
advocacy career on the same premise through the Joyful Heart Foundation
and the campaign against the rape kit backlog. When the culture caught
up to the show's stated ethic in the fall of 2017, SVU became the
natural test case for a question that matters well beyond one
procedural: what does a believability movement do to the stories a
believability show tells about accusations that turn out to be wrong?

Two predictions circulated, and they pointed in opposite directions. The
first came out of the backlash literature. If popular culture absorbs
the anxieties of the moment, and if the loudest anxiety after October
2017 was the specter of the falsely accused man, then crime drama should
have leaned into wrongful accusation plots: more of them, more
sympathetic to the accused, more skeptical of accusers (Boyle 2019;
Banet-Weiser 2018; Manne 2017). The second prediction arrived three
years later. After the murder of George Floyd in 2020, the police
procedural came under sustained criticism as copaganda, and advocacy
reports urged writers' rooms to depict less misconduct, more
accountability, and gentler cops (Color of Change and USC Annenberg
Norman Lear Center 2020). On that account, whatever SVU did with its
accusation plots, its detectives should at least have softened.

Neither prediction survives contact with the coded record. Across the
three eras SVU cut its false-suspect episode rate by more than a third,
and the density of wrongly accused persons per episode fell by half. The
cases that survived got harsher, with severe or life-altering outcomes
rising above eight in ten during Seasons 19-21. The detectives,
meanwhile, did the opposite of softening: depicted threats and coercion
against innocent suspects, which had drifted to a floor in the
mid-2010s, rebounded after 2020 to early-series levels. Section 4
reports the series and the tests. The complication is timing. Seasons 19
and 20, written and aired at the height of the movement, contain false
suspects at rates equal to or above the seasons just before them; the
collapse arrives in Season 21, two years into the era. Season 21 is also
Warren Leight's first season back as showrunner, and its production
ended early in the COVID shutdown. Movement, showrunner, pandemic: any
of the three could be holding the pen, and season-block data cannot
separate them.

What the design can establish is still worth having. The configuration
of the show changed across the \#MeToo boundary, in frequency, in
severity, in whom the scripts blame, and in how the police behave, and
both of the standard narratives fail regardless of which driver
dominates. I will argue that the pattern is best read as risk management
by deletion, the cheap weekly version of the false accusation plot cut
and the catastrophic version retained, but two rival mechanisms, a
showrunner's taste and a measurement artifact I call reform-censoring,
get a full hearing in Section 5.

The analysis stands on a census rather than a sample. Every aired
episode was coded, and after a source audit excluded nine episodes whose
transcripts belong to other episodes, the analysis set is 567 episodes
and 521 falsely accused person-episode records. A companion paper
(Orden, in preparation) reports the descriptive anatomy of the
phenomenon: the overall rates, the severity distribution, the death
toll, the apology deficit. This paper takes those measures across time.

\subsection{2. Background}\label{background}

\subsubsection{2.1 SVU and the scholarship of the special
victim}\label{svu-and-the-scholarship-of-the-special-victim}

SVU has attracted academic attention almost since its premiere, most of
it focused on victims rather than suspects. Britto and colleagues (2007)
found the show's victim population skewed heavily toward young white
women relative to victimization statistics. Cuklanz and Moorti (2006)
read the early seasons as a hybrid text, feminist in its insistence that
rape is grave and disbelief is harmful, yet repeatedly drawn to plots
that recentered male authority and punished maternal failure. Their
later book-length treatment traces how the show's politics moved with
its decades (Moorti and Cuklanz 2017). The genre context runs deeper:
Cuklanz (2000) shows prime-time rape narratives remaking themselves
decade by decade, with the sensitive detective becoming the proof of
enlightenment, and Projansky (2001) documents how thoroughly rape plots
structure postfeminist television. The broader procedural literature
establishes the baseline: prime time crime drama exaggerates clearance
rates, personalizes justice, and treats police process as trustworthy
machinery (Eschholz, Mallard, and Flynn 2004; Surette 2015).

SVU matters for this literature out of proportion to its slot on the
schedule. It is the longest-running primetime live-action series in
American television, syndication has never really taken it off the air,
and its detectives are among the most recognizable fictional police in
the world. A representational habit on SVU is a habit rehearsed for tens
of millions of viewers over a quarter century. The show is also unusual
in carrying an explicit ideological signature: where the parent Law \&
Order sold procedural neutrality, SVU sells advocacy, an hour built
around the premise that these particular victims have historically been
doubted and deserve better.

What the scholarship has largely skipped is the falsely accused person
as a recurring character type. The omission is understandable, since no
one had counted them. It turns out there are 521 qualifying
person-episodes on SVU alone, nearly one per episode across the full
run, and their treatment varies by era in ways that track the show's
cultural surroundings. That variation is the subject here.

\subsubsection{2.2 \#MeToo as a problem for a believability
show}\label{metoo-as-a-problem-for-a-believability-show}

The \#MeToo movement put a specific proposition into general
circulation: accusations of sexual violence are overwhelmingly true, and
the reflex to doubt them is itself a mechanism of impunity. Empirically
the movement stood on firm ground; the research consensus places
knowingly false reports of sexual assault in the low single digits to
roughly ten percent of cases (Lisak et al.~2010). Culturally, the
proposition collided with a backlash that organized itself around the
mirror-image figure, the innocent man ruined by a lie. Manne (2017) gave
the reflex a name, himpathy, and Boyle (2019) documented how quickly due
process talk became the respectable costume of accuser skepticism.

A show like SVU sat directly on this fault line. Its brand was belief.
Its plot engine, though, ran substantially on doubt: red herrings, wrong
arrests, recanted identifications, the mid-episode pivot from one
suspect to the real one. Every such beat dramatizes a case where
believing the initial accusation would have convicted the wrong person.
Before 2017 this tension was invisible because no one was counting.
After 2017, a writers' room producing a false accusation plot was making
a legible political choice, and the record examined below suggests the
room knew it.

The backlash discourse gave that choice a public price. Within a year of
the Weinstein reporting, the falsely accused man had become the
organizing figure of \#HimToo commentary, of the Kavanaugh confirmation
fight, and of a due process argument prosecuted largely through
anecdote. Any episode in which an accuser turned out to be lying or
mistaken was raw material for that argument, clippable and shareable
regardless of the episode's own framing. Fictional false accusations had
become citable, and a show whose stars campaigned for survivors had
reasons to think hard before supplying the citations.

\subsubsection{2.3 The copaganda debate after
2020}\label{the-copaganda-debate-after-2020}

The second cultural shock arrived in the summer of 2020. The procedural
genre, long criticized in passing, became the object of a focused
campaign. Color of Change's Normalizing Injustice report (2020) had
already documented how rarely scripted crime shows depict police
misconduct as wrongful, and it carried an explicit prescription: less
glorified aggression, more accountability. What actually happened on air
was varied rather than uniform; Bernabo (2022) documents a range of
narrative strategies across the 2020-21 broadcast responses, not a
single reform posture.

The reform prescription is unusually easy to operationalize, so it can
be tested directly. If the procedural reformed after 2020, three things
should appear in coded data: depicted threats and coercion against
suspects should fall, apologies and other accountability beats should
rise, and the misconduct that remains should attach to characters the
narrative marks as wrong. The first two are directly measurable here.
The third requires framing analysis beyond this instrument, and I return
to that gap in the discussion, but a genre that reformed on the first
two margins would at minimum stop escalating them.

SVU makes a good test because its misconduct is measurable at the person
level. Each coded person carries a field recording whether detectives
threatened, deceived, insulted, or otherwise coerced them, with the
supporting quote preserved. Since every person in the dataset is
innocent, the measure captures the show's willingness to depict
aggressive policing aimed at the people reform discourse worries about
most: those the aggression cannot, even within the fiction's own logic,
be said to have deserved.

\subsubsection{2.4 Cultivation and the rare
event}\label{cultivation-and-the-rare-event}

Cultivation theory holds that heavy television exposure shapes viewers'
estimates of social reality, especially for phenomena the viewer rarely
encounters firsthand (Gerbner et al.~2002; Morgan and Shanahan 2010).
False accusations of sexual violence are such a phenomenon. Almost no
viewer has independent knowledge of their frequency; the fiction is the
data. A show that airs a false suspect in two episodes out of three
makes one base rate available to its audience, and a show that thins
those suspects out while making their outcomes harsher makes a different
one available. The longitudinal design used here is, in effect, a
measurement of the changing cultivation dose.

\subsection{3. Method}\label{method}

Full instrument detail, including the complete codebook and the
descriptive results, appears in the companion paper (Orden, in
preparation); the codebook itself accompanies this submission as
supplementary material. The essentials follow.

\textbf{Data and source quality.} The corpus is a census of all 576 SVU
episodes, Seasons 1 through 27, Season 27 partial at three episodes
because the season was still in progress in January 2026. Transcripts
were obtained from subslikescript.com and stored in a spreadsheet, one
row per episode. A ground-truth audit of that source, run after coding,
found nine episodes whose cells contain a duplicate of another episode's
script, a corruption the coding inherited silently. The nine (all in
Seasons 1-18) are excluded from every episode denominator, and the 15
person-records they generated are removed. Four further episodes had
empty transcripts and one had corrupted encoding, for 14 defective
source episodes in all; the four empty ones contribute no person records
and slightly deflate episode rates in their seasons. The audit also
found that 145 of 576 transcripts are truncated at exactly 32,767
characters, the Excel cell limit, and the truncation is lopsided: 138 of
the 145 fall in Seasons 1-18 (33.7\% of that era) against 2 in Seasons
19-21 and 5 in Seasons 22-27. Truncation cuts endings, where
late-episode harms and apologies live, so it can only depress the
early-era counts; its effect on this paper's central finding is
conservative. One artifact the audit ruled out directly: transcript
length by season shows no downward era trend (era means of 5,682, 5,381,
and 5,662 words per episode), so the retreat documented below is not an
artifact of thinner transcripts.

\textbf{Coding.} Each transcript was coded by a large language model
(claude-sonnet-4-5, Anthropic) via batch API under a fixed codebook with
controlled vocabularies, a design that sits inside the emerging
LLM-annotation methods literature and its validation norms (Gilardi,
Alizadeh, and Kubli 2023; Törnberg 2023; Ziems et al.~2024), with
classical content-analytic standards as the eventual benchmark
(Krippendorff 2019; Neuendorf 2017). The inclusion rule admits a person
only if three conditions hold: proven or strongly implied innocent of
the accused offense; subjected to public exposure or a physical
consequence; and harmed by the accusation itself rather than by
unrelated events. Each record carries the accusation's origin under a
controlled vocabulary (victim identification, witness misidentification,
squad inference, coerced statement, database error, fabrication), the
exposure channel, police conduct toward the person (none, verbal threat,
coercive tactic, insult, multiple), whether police apologized, and
consequence severity on a four-point scale: 1, private harm; 2, public
exposure without formal sanction; 3, material sanction such as arrest or
firing; 4, life-altering outcome or death. The raw file contains 541
person-records. Five coded actually guilty are excluded, as are the 15
phantom records from the corrupted episodes, leaving n = 521
person-episode records, roughly 513 unique persons once recurring
characters are deduplicated. The 43 partially involved records are
retained: such persons are innocent of the accusation, whatever lesser
wrong they committed. The model rates its own coding confidence high on
97\% of records; I give that number no evidentiary weight pending the
planned human validation pass. All quoted dialogue below was checked
verbatim against the source transcripts, which carry no speaker labels;
named attributions follow scene context.

\textbf{Eras.} The analysis files carry no air dates (the raw coding
output has them for only 121 of 576 episodes, mostly early seasons), so
eras are operationalized by production season: Seasons 1-18 (1999
through spring 2017) as the pre-\#MeToo era, Seasons 19-21 (fall 2017
through 2020) as the \#MeToo era, and Seasons 22-27 as the post-2020
era. Season blocks are a coarse proxy, conservative in a useful
direction: Season 19 premiered days before the Weinstein reporting, so
its earliest episodes were written in a pre-\#MeToo world, and any
contamination biases the era contrast toward zero. An air-date join from
external sources would permit true event-time analysis and showrunner
segmentation; that is planned follow-up work.

\textbf{Sensitivity.} Every episode-level rate is computed under two
treatments of the flag's Maybe value, counting Maybe episodes as
negative (Y-only) and as positive (Y+Maybe). At the era level the
treatments differ by 3.0 points in the pre-\#MeToo block, 2.9 in Seasons
19-21, and 1.0 in Seasons 22-27; every substantive conclusion holds
under both, and the chi-square results are reported both ways.

\textbf{Analysis and inference.} The dataset is a census, so a word on
why tests appear at all. I treat each season as one draw from the
writers'-room process that could have generated other seasons; on that
superpopulation view, the tests ask whether the era contrasts exceed the
noise such a process produces. Persons cluster within episodes, so the
person-level tests are indicative and the episode-level tests are
primary. Era comparisons use Pearson chi-square with two degrees of
freedom; the smallest expected cell across all reported tests is 14.7,
well above the standard minimum of five. Four tests were planned. The
two episode-level results survive any multiple-comparison correction;
the severity contrast (p = .018) would not survive a Bonferroni
correction across the four (.018 x 4 = .07), so I treat it as
suggestive; the coercion test is nonsignificant outright. Within-era
season cells shrink sharply after Season 14 (3 to 18 persons per
season), so single-season values are reported descriptively and
inference rests on the era and three-season-block aggregates. Era person
counts for Seasons 19-21 and 22-27 are 39 and 55; I treat percentage
swings under roughly ten points in those columns as noise.

\subsection{4. Results}\label{results}

\subsubsection{4.1 The retreat}\label{the-retreat}

The per-season series in Table 1 is the primary evidence. The
pre-\#MeToo era is peaked rather than flat: the false-suspect rate runs
above 80\% of episodes in Seasons 8 through 10 and has already eased to
between 47.8\% and 54.2\% across Seasons 14-18. Then come the two
seasons written under \#MeToo's full glare, and they do not flinch:
Season 19 posts 54.2\% and Season 20 posts 62.5\%, at or above the
plateau of the five seasons before them. The break lands at Season 21,
which collapses to 30.0\%, the series floor. Season 21 is the
pandemic-shortened 2019-20 season and Warren Leight's first as returning
showrunner. Every season since has stayed below the old plateau except
Season 26, which rebounds to 54.5\% on 22 episodes, a single-season
wobble worth watching.

Pooled into eras, the retreat is large and statistically secure (Table
2). A false suspect appears in 65.3\% of pre-\#MeToo episodes (262 of
401), 50.0\% of \#MeToo-era episodes (34 of 68), and 41.8\% of post-2020
episodes (41 of 98); the era association is strong (chi-square(2, N =
567) = 20.89, exact p = .000029) and strengthens under the inclusive
Y+Maybe denominator (68.3\% to 52.9\% to 42.9\%; chi-square = 24.45,
exact p = .000005). But the pooling should not be allowed to relocate
the break: the era means straddle Season 21, and it is Season 21, not
the S18/S19 boundary, where the step down occurs.

The person-level measure fell further and earlier. Seasons 1-18 average
1.06 falsely accused persons per episode against 0.57 in Seasons 19-21
and 0.56 since, but the per-season column of Table 1 shows the density
already sliding through the mid-2010s, from crowded early years (1.53
per episode in Season 2, 1.89 in Season 9) to the 0.6-0.8 band across
Seasons 14-20. What the \#MeToo-era seasons add is the floor: six coded
persons across twenty hours of Season 21, 0.30 per episode.

Read together, the two measures describe a change in editorial diet
rather than a blip. For eighteen seasons the wrongly accused person was
a stock ingredient, present more often than absent. Somewhere between
the fall of 2019 and the fall of 2020, the ingredient became occasional,
and it has stayed occasional for six years. Only the coercion series
(Section 4.4) moves earlier.

\subsubsection{4.2 The severity paradox}\label{the-severity-paradox}

The retreat did not make life easier for the falsely accused characters
who remained. Among coded persons, the share suffering a material
sanction or a life-altering outcome (severity 3-4) rises from 60.0\%
before \#MeToo (256 of 427) to 82.1\% during it (32 of 39), then settles
at 67.3\% after 2020 (37 of 55), a level the paper's own noise rule
treats as indistinguishable from baseline (Table 3). The era association
is suggestive rather than secure (chi-square(2, n = 521) = 8.06, p =
.018, uncorrected). Mean severity follows the same arc, 2.67 to 3.05 to
2.84 on a four-point scale (the median is 3 in all three eras), and the
share at severity 4 alone, the level reserved for death, permanent ruin,
or severe injury, jumps from 20.1\% to 30.8\% before returning to
21.8\%. The precise claim, then, is bounded: severity rose sharply
during the \#MeToo-era seasons and has settled no milder than it ever
was.

The \#MeToo-era cases carry the pattern on their faces. Greg Harvey
(S19E07) is murdered while under false suspicion. Ricardo Torres
(S21E06) was coerced into false testimony at fifteen, spent fifteen
years in prison, contracted HIV there, and dies of AIDS shortly after
release. These are the plots SVU chose to keep while cutting the
category in half.

The post-2020 seasons relax the conditional severity only partway, and
they keep producing bodies. Lonnie Liston (S22E10) is tortured and
murdered by vigilantes after the squad reopens him as a suspect. Season
26 alone has three: vigilantes attack Eddie Soto (S26E13), Father
Alberto is murdered (S26E19), and Anthony Fierro dies by suicide
(S26E22).

One adjacent series can be dispatched in a sentence. The police apology
rate barely moves across eras, 7.5\% of persons, then 5.1\%, then 9.1\%,
on era cells far too small to distinguish; whatever changed after 2017,
contrition was no part of it.

\subsubsection{4.3 The origin shift}\label{the-origin-shift}

Table 5 breaks down who first points the finger. In every era the squad
itself remains the leading source: detectives' own inference produces
56.7\% of false accusations before \#MeToo, 41.0\% during it, and 47.3\%
after 2020. The show's single most consistent finding, across 27
seasons, is that the wrongly accused person is most often accused by the
police.

The movement happens in two smaller categories, and it is directional.
Accusations originating in the victim's own identification of the wrong
person roughly double their share at the \#MeToo boundary, from 15.5\%
(66 of 427) to 30.8\% (12 of 39), and hold at 29.1\% (16 of 55) after
2020. Fabrication, the knowingly false claim, moves the opposite way on
a delay: 13.1\% before \#MeToo, 17.9\% during it, then 5.5\% (3 of 55)
after 2020. The two movements deserve different weights. The
victim-identification doubling clears the paper's ten-point noise rule
and persists in the larger post-2020 cell; it is the finding. The
fabrication decline, at 7.6 points from baseline, sits inside the noise
band for a 55-person column, so I report it as consistent with the same
editorial turn and build nothing on it. The remaining categories hold
steady or move within noise.

The secure part of the shape is this: after \#MeToo, when SVU airs a
false accusation, it increasingly locates the error in a sincere victim
who is honestly mistaken. The fabrication point estimates say the liar
is leaving the repertoire too, though the noise band withholds judgment
on that half. Mehcad Carter (S15E03) prefigures the type: a Black
teenager identified from photographs by two separate victims, named
publicly, and shot dead by a woman who believed she was defending
herself from the rapist the identification had made him. The accusers in
that story are neither villains nor liars. They are wrong, and the later
eras return to that configuration again and again.

\subsubsection{4.4 A U-shape in depicted
coercion}\label{a-u-shape-in-depicted-coercion}

The police conduct series is the most surprising in the dataset, and its
statistical caveat comes first: the pooled era test does not reach
significance (chi-square(2, n = 521) = 5.30, p = .071), and the claim
here rests on the three-season-block series, a shape the pooled test
cannot see because the decline begins years before the era boundary and
pooling buries it inside the pre-\#MeToo block.

Table 4 shows the blocks. From Seasons 1 through 12 the threat rate is
stable and high, between 57.7\% and 60.3\% of persons per block. The dip
begins at Season 15 (26.7\%) and bottoms out across Seasons 16-18 at
27.7\% (13 of 47), with Season 16 the gentlest season in show history at
14.3\%. Recovery is only partial in the \#MeToo block (35.9\%). Then,
after 2020, the rate snaps back: 56.7\% in Seasons 22-24 and 60.0\% in
Seasons 25-27, descriptively level with the early-series plateau. The
five lowest single seasons are 16, 24, 15, 19, and 18; four of the five
sit in the mid-2010s trough, and the fifth, Season 24 at 23.1\%, is an
isolated low inside an otherwise coercive rebound era. Within that
rebound era the single-season values whipsaw between 87.5\% and 23.1\%
on cells of eight to fourteen persons (see the note to Table 4), and the
S25-27 block leans on a Season 27 that is three episodes and three
persons old. The blocks smooth this to 56.7\% and 60.0\%, and it is the
blocks I trust.

The rebound-era episodes supply the texture. In S22E10, detectives tell
Lonnie Liston, a Black exoneree under fresh suspicion, ``Better hope she
shows up. Alive.'' Later in the same hour one tells him, ``Oh, you don't
have a job anymore. And if I were you, I wouldn't go out there,'' as
protestors chant outside. Liston is then surveilled, threatened with a
parole violation, and finally tortured and murdered by vigilantes. In
S23E07, detectives tell Country Jones they are tearing apart his
restaurant and his car, and that a lie to police is enough to revoke the
last three months of his parole. These scenes would be unremarkable in
Season 4. They aired in 2020 and 2021.

So the record shows SVU's gentlest policing occurring in roughly
2013-2017, before either \#MeToo or the Floyd protests, followed by a
post-2020 return to early-series levels of on-screen coercion against
the innocent.

\subsubsection{4.5 Media exposure and the remaining
distributions}\label{media-exposure-and-the-remaining-distributions}

One exposure channel moves sharply with the eras. The share of falsely
accused persons exposed through the media jumps from 15.2\% before
\#MeToo to 33.3\% (13 of 39) in Seasons 19-21, then falls back to 14.5\%
after 2020; counting the two online-exposure records, the combined share
runs 15.5\% to 33.3\% to 16.4\% (Table 6). During the years when the
culture was arguing about trial by media, one falsely accused SVU
character in three was exposed through the press, double the show's
baseline. The companion paper shows media exposure is the most dangerous
channel in the dataset, with a severity-4 rate 5.4 times the police-only
rate, so the spike compounds the era's conditional severity.

The rest of the distributions move little. The accusation type
concentrates, with rape rising from 57.8\% of accusations before \#MeToo
to 71.8\% and then 74.5\% as the peripheral assault and vaguely coded
sex-crime cases thin out alongside the red herrings; workplace, family,
and school exposure hold roughly steady; and the squad remains the party
that tells the world in the overwhelming majority of cases throughout.

\subsection{5. Discussion}\label{discussion}

\subsubsection{5.1 Deletion, and two
rivals}\label{deletion-and-two-rivals}

The cleanest summary of Sections 4.1 and 4.2 is that SVU stopped telling
the cheap version of the false accusation story and kept the expensive
one. The pre-\#MeToo show ran on procedural churn: most episodes cycled
through at least one wrong suspect, and most of those suspects suffered
exposure or arrest and then vanished from the story, cleared and
forgotten. Those beats became legible, after October 2017, as small
weekly arguments that accusations are unreliable. A writers' room that
wanted no part of that argument had an easy remedy: delete the casual
false suspect, keep the wrongful accusation as tragedy, the story in
which the show itself stands with the destroyed innocent. The
per-episode dose halves; conditional severity climbs 22 points during
the era of maximum scrutiny. On the deletion reading, that pair is one
editorial decision.

The timing evidence supports the deletion reading less cleanly than a
movement story would like, and the claim has to carry the rivals. The
break is Season 21, not Season 19: the two seasons produced at the
height of the movement kept staging false suspects at the old late-era
rate, and the collapse coincides with a returning showrunner and a
pandemic shutdown. Leight's tenure is a fully live explanation, and so
is COVID-era production, though per-episode denominators strip out the
mechanical effect of shortened seasons. What the timing does rule out is
reading the retreat as an old trend completing itself: Seasons 14-18 had
plateaued, Seasons 19-20 sat at and above that plateau, and the step
down is abrupt, deep, and durable, holding through subsequent turnover
in the room. The deletion outlasted its occasion, whoever signed it.

A second rival lives in the instrument rather than the room, and it
needs a name: reform-censoring. The inclusion rule requires exposure or
physical consequence before a wrongly suspected person becomes a record.
A show that kept suspecting innocent people but stopped exposing them, a
show that had actually reformed, would mechanically produce fewer
records at higher conditional severity, exactly the observed pattern.
The episode-level flag discriminates, because it does not depend on the
person-level inclusion rule. Within flagged episodes, person-records per
episode fall from 1.62 before \#MeToo to 1.15 in Seasons 19-21,
recovering to 1.34, and flagged episodes with zero qualifying persons
rise from 5.0\% to 11.8\% before easing to 7.3\%. That is a real
censoring signal in Seasons 19-21, resting on four of 34 episodes, and
it may account for part of the era's severity spike. It cannot account
for the retreat: the headline decline is carried by the episode-level
flag itself, which fell from 65.3\% to 41.8\% with no help from the
inclusion rule.

One bias runs in the paper's favor. A third of pre-era transcripts are
truncated at the Excel cell limit, against a handful later, and
truncation deletes endings, where harms and apologies concentrate;
early-era harm is if anything undercounted, and the documented decline
is conservative.

\subsubsection{5.2 The origin shift and the licensing of
error}\label{the-origin-shift-and-the-licensing-of-error}

The shift documented in Section 4.3 shows the same instinct operating
inside the surviving plots. A false accusation story needs someone to be
wrong. The available candidates are the accuser, a witness, a database,
or the police, and the choice among them is ideological whether or not
anyone in the room says so. Before \#MeToo, SVU spread the error around,
with a meaningful minority of plots resting on a fabricating accuser.
The shift that clears the noise rule sits on the other side of the
ledger: the sincere but mistaken victim carries twice her former share
after 2017 and holds it after 2020. The fabricator's slide to 5.5\%
points the same way but stays inside the noise band, so the argument
rests on the doubling alone, and the doubling is enough. The
post-\#MeToo show has, in effect, found a formulation under which it can
keep telling wrongful accusation stories without contradicting its
brand: on this show, accusers err. Belief in the victim's sincerity is
preserved, and the cost lands elsewhere, because someone still gets
destroyed, and the destroyed man of the later eras suffers more than his
predecessor did.

Note what did not shift. Squad inference remains the leading origin in
every era. The show never stopped telling the story in which its own
heroes point the first finger at an innocent person; it simply never
framed that story as being about the pointing. The institutional
self-scrutiny that \#MeToo might have prompted, an interrogation of the
detectives' own accusatory power, is absent from the trend lines.

The dataset lacks systematic demographic coding, but one pattern in the
qualitative record is hard to miss. Many of the catastrophic
misidentification cases across the later pre-\#MeToo years and the
rebound era involve men of color: Mehcad Carter, Jerome Jones (S16E13),
Carlos Hernandez (S21E06), Lonnie Liston. The sincere-mistake plot,
whatever its intentions toward accusers, keeps finding the same bodies
to land on. Coding race and gender of the accused systematically is the
most important extension of this instrument, and I make no quantitative
claim here beyond the observation that the era's exemplary victims of
error are conspicuously patterned.

\subsubsection{5.3 The rebound against the reform
narrative}\label{the-rebound-against-the-reform-narrative}

Section 4.4 is the finding I would most like a second dataset to check,
because it inverts the received account of the post-2020 procedural. The
reform campaign predicted a softening: fewer glorified rough
interrogations, more visible accountability (Color of Change and USC
Annenberg Norman Lear Center 2020), and Bernabo (2022) records how
varied the networks' actual responses were. On SVU the softening
happened, but it happened in roughly 2013 through 2017, bottoming at a
Season 16 in which one coded person in seven faced any threat. After
2020 the show returned to depicting coercion against the innocent at the
rate of the early Stabler years, and the apology rate stayed flat
throughout.

The dip itself wants explaining as much as the rebound does. The gentle
years, Seasons 15 through 18, were written as Ferguson and its aftermath
were making police violence a national argument, and long after the
show's own center of gravity had moved from Stabler's fists to Benson's
empathy. I read the trough as SVU's actual reform era, unannounced and
probably unplanned as such, arriving years before the discourse that
supposedly demanded one. By the time the culture asked the procedural to
soften in 2020, this procedural had already softened and was on its way
back.

Two interpretations of the rebound fit the coding, and the
transcript-only instrument cannot separate them. The first is
recommitment: the show reverted to aggressive-interrogation drama
because it is reliable television. The second is critique: the later
seasons may stage coercion in order to indict it, and the Liston episode
is a plausible example, since the vigilante murder of an exoneree reads
as an indictment of suspicion itself. The distinction matters for
authorial intent; it matters less for cultivation. Under either reading,
a viewer of Seasons 22-27 watches innocent people threatened and coerced
by police at roughly double the rate a viewer of Seasons 16-18 did
(58.2\% of coded persons against 27.7\%; about 1.7 times per episode
aired, since the later seasons also carry fewer accused persons per
hour). The exposure is the dose, and the dose went up after the summer
that was supposed to bring it down.

\subsubsection{5.4 Who the viewer learns to
trust}\label{who-the-viewer-learns-to-trust}

Cultivation theory earns its place here because false accusation
frequency is a base-rate belief with real consequences, held almost
entirely on fictional evidence. Jurors, employers, and family members
act on private estimates of how often accusations are wrong, and almost
none derive them from case files. The pre-\#MeToo SVU offered a weekly
exemplar of accusation landing on the wrong person: two episodes in
three, roughly one innocent per hour, for eighteen years. Comparing that
to the documented record takes care, because a false suspect is not a
false report: SVU's wrongly accused are mostly the squad's own
candidates, and the 2-to-10-percent consensus on knowingly false reports
(Lisak et al.~2010) speaks only to the fabrication slice of the
repertoire, which ran at 13.1\% of coded persons before \#MeToo and at a
point estimate of 5.5\% after 2020, a decline consistent with the era
story though inside the noise band for a 55-person cell. What the old
show oversupplied, relative to anything in the record, was the spectacle
of wrongful suspicion itself, and it made that spectacle available
weekly to anyone forming a base-rate belief. Whether viewers formed the
belief is a reception question this design cannot answer, and the
audience evidence warns against assuming they did: Hust et al.~(2015)
found crime-drama franchise exposure associated with lower rape-myth
acceptance and better consent negotiation among college students.
Message is not uptake.

The current message system is different. Wrongful accusation is now rare
on screen, which brings the fictional base rate closer to the empirical
one, but the exemplars are vivid in the way that dominates memory: the
perp walk, the prison assault, the vigilante killing, the sixteen lost
years. The sincere, mistaken accuser drives twice her old share of the
wreckage; the lying accuser appears headed out of the repertoire, though
that decline sits within noise; and the police who generate most of the
false accusations neither soften nor apologize. Whether that
configuration serves accusers or the accused is genuinely unclear to me.
What is clear is that it was assembled measurably, one editorial
decision at a time, on a calendar that tracks the culture's.

\subsection{6. Limitations}\label{limitations}

Everything above rests on one show. SVU is the genre's longest-running
and most belief-branded title: the right single case for this question,
and a poor basis for claims about the genre at large. Era assignment is
by season block because the analysis files carry no air dates, so
episodes are dated by production season, and the Season 19 boundary
mixes a handful of pre-Weinstein scripts into the \#MeToo era, biasing
those contrasts toward zero. The two later eras contain 39 and 55
persons, so their percentage estimates are coarse; I have flagged which
movements clear the noise threshold, the pooled coercion test is
nonsignificant, and the U-shape claim rests on the block series and is
exploratory.

The coder is a large language model working from transcripts alone, and
it carries two distinct risks. First, the measured instruction
violations on checkable fields (prohibited guilty records returned,
out-of-vocabulary values, formatting violations) imply a nonzero error
rate on uncheckable fields, so the planned human validation pass is a
requirement for confidence, not a formality. Second, and specific to
this paper: the coder saw season and episode numbers and carries
pretraining knowledge of both SVU and \#MeToo, so any coding drift could
itself be era-correlated, which is the axis this paper measures. The
planned check is to re-code a stratified subsample with season and
episode metadata stripped. Transcript-only coding is also blind to
performance, camera, and score, which is the information needed to
separate depicted coercion from endorsed coercion.

The source defects cut in known directions. The nine wrong-transcript
episodes are excluded outright; the transcript-length check shows no era
trend in transcript thickness; and the heavy pre-era truncation biases
the retreat estimate toward zero, since the era with a third of its
transcripts cut short is the era whose harms are undercounted. Four
episodes lack transcripts entirely, one is encoding-damaged, and Season
27 is three episodes old. None of these problems seems likely to
manufacture a 23-point drop in the false-suspect rate or a doubling of
the victim-identification share, but each one argues for holding the
point estimates loosely.

\subsection{7. After the Retreat}\label{after-the-retreat}

Three extensions would harden or overturn these claims, and I name them
in order of value. First, the same instrument applied to other
long-running procedurals would establish whether the
retreat-and-intensify pattern is SVU's alone or the genre's general
answer to \#MeToo; preliminary work on the parent Law \& Order suggests
the franchises diverged, which would make SVU's pattern a choice rather
than an industry tide. Second, an air-date and showrunner join would
directly test the Season 21 rivals that the season-block design leaves
standing. Third, a framing-sensitive recoding of the police conduct
quotes, asking whether anyone on screen objects to the coercion, would
settle the endorsement question that Section 5.3 leaves open.

The larger lesson is about how a long-running institution edits itself
under pressure. SVU did not argue with \#MeToo, and it did not visibly
reform after 2020. It curated. It deleted the cheap version of a newly
dangerous plot, kept the expensive version, relocated the blame inside
it, and let its detectives revert to form once the spotlight moved;
whether the curator was the movement, the showrunner, or the pandemic,
the curation is in the record, and each move is invisible at the level
of any single episode and unmistakable in the aggregate. Counting is the
only criticism that catches this kind of authorship. The record says the
show chose retreat. The few it still accuses pay no less than they ever
did, and for three seasons they paid more.

\subsection{Acknowledgments}\label{acknowledgments}

Claude (Anthropic) served two roles in this project, and I want to state
both plainly. First, as measurement instrument: a Claude model
(claude-sonnet-4-5) coded all 576 episode transcripts under the fixed
codebook described in Section 3, a use I treat as I would any automated
coder, pending the human validation described in the limitations.
Second, as drafting assistant: Claude produced an initial draft of this
text under my direction, which I revised and verified against the
underlying dataset. The same model family thus coded the data and
drafted the text, a circularity that is one more reason the planned
validation cross-check uses a non-Anthropic model. All analytic
decisions, interpretations, and errors are mine, and I take full
responsibility for the content.

\subsection{References}\label{references}

Banet-Weiser, S. 2018. \emph{Empowered: Popular Feminism and Popular
Misogyny}. Durham, NC: Duke University Press.

Bernabo, L. 2022. ``Copaganda and Post-Floyd TVPD: Broadcast
Television's Response to Policing in 2020.'' \emph{Journal of
Communication} 72 (4): 488-496.

Boyle, K. 2019. \emph{\#MeToo, Weinstein and Feminism}. Cham: Palgrave
Macmillan.

Britto, S., T. Hughes, K. Saltzman, and C. Stroh. 2007. ``Does `Special'
Mean Young, White and Female? Deconstructing the Meaning of `Special' in
Law \& Order: Special Victims Unit.'' \emph{Journal of Criminal Justice
and Popular Culture} 14 (1): 39-57.

Color of Change and USC Annenberg Norman Lear Center. 2020.
\emph{Normalizing Injustice: The Dangerous Misrepresentations That
Define Television's Scripted Crime Genre}. Los Angeles: USC Annenberg.

Cuklanz, L. M. 2000. \emph{Rape on Prime Time: Television, Masculinity,
and Sexual Violence}. Philadelphia: University of Pennsylvania Press.

Cuklanz, L. M., and S. Moorti. 2006. ``Television's `New' Feminism:
Prime-Time Representations of Women and Victimization.'' \emph{Critical
Studies in Media Communication} 23 (4): 302-321.

Eschholz, S., M. Mallard, and S. Flynn. 2004. ``Images of Prime Time
Justice: A Content Analysis of NYPD Blue and Law \& Order.''
\emph{Journal of Criminal Justice and Popular Culture} 10 (3): 161-180.

Gerbner, G., L. Gross, M. Morgan, N. Signorielli, and J. Shanahan. 2002.
``Growing Up with Television: Cultivation Processes.'' In \emph{Media
Effects: Advances in Theory and Research}, 2nd ed., edited by J. Bryant
and D. Zillmann, 43-67. Mahwah, NJ: Lawrence Erlbaum.

Gilardi, F., M. Alizadeh, and M. Kubli. 2023. ``ChatGPT Outperforms
Crowd Workers for Text-Annotation Tasks.'' \emph{Proceedings of the
National Academy of Sciences} 120 (30): e2305016120.

Hust, S. J. T., et al.~2015. ``Law \& Order, CSI, and NCIS: The
Association Between Exposure to Crime Drama Franchises, Rape Myth
Acceptance, and Sexual Consent Negotiation Among College Students.''
\emph{Journal of Health Communication} 20 (12): 1369-1381.

Krippendorff, K. 2019. \emph{Content Analysis: An Introduction to Its
Methodology}. 4th ed.~Thousand Oaks, CA: Sage.

Lisak, D., L. Gardinier, S. C. Nicksa, and A. M. Cote. 2010. ``False
Allegations of Sexual Assault: An Analysis of Ten Years of Reported
Cases.'' \emph{Violence Against Women} 16 (12): 1318-1334.

Manne, K. 2017. \emph{Down Girl: The Logic of Misogyny}. New York:
Oxford University Press.

Moorti, S., and L. Cuklanz. 2017. \emph{All-American TV Crime Drama:
Feminism and Identity Politics in Law and Order: Special Victims Unit}.
London: I.B. Tauris.

Morgan, M., and J. Shanahan. 2010. ``The State of Cultivation.''
\emph{Journal of Broadcasting \& Electronic Media} 54 (2): 337-355.

Neuendorf, K. A. 2017. \emph{The Content Analysis Guidebook}. 2nd
ed.~Thousand Oaks, CA: Sage.

Orden, T. In preparation. ``Collateral Damage as Narrative Convention:
False Accusation and Its Costs Across 27 Seasons of Law \& Order:
Special Victims Unit.'' Manuscript.

Projansky, S. 2001. \emph{Watching Rape: Film and Television in
Postfeminist Culture}. New York: NYU Press.

Surette, R. 2015. \emph{Media, Crime, and Criminal Justice: Images,
Realities, and Policies}. 5th ed.~Stamford, CT: Cengage.

Törnberg, P. 2023. ``ChatGPT-4 Outperforms Experts and Crowd Workers in
Annotating Political Twitter Messages with Zero-Shot Learning.''
arXiv:2304.06588.

Ziems, C., W. Held, O. Shaikh, J. Chen, Z. Zhang, and D. Yang. 2024.
``Can Large Language Models Transform Computational Social Science?''
\emph{Computational Linguistics} 50 (1): 237-291.

\begin{center}\rule{0.5\linewidth}{0.5pt}\end{center}

\subsection{Tables}\label{tables}

\textbf{Table 1. False-suspect rate and falsely accused persons per
episode, by season (primary series).}

{\def\LTcaptype{none} % do not increment counter
\begin{longtable}[]{@{}
  >{\raggedright\arraybackslash}p{(\linewidth - 8\tabcolsep) * \real{0.2000}}
  >{\raggedright\arraybackslash}p{(\linewidth - 8\tabcolsep) * \real{0.2000}}
  >{\raggedright\arraybackslash}p{(\linewidth - 8\tabcolsep) * \real{0.2000}}
  >{\raggedright\arraybackslash}p{(\linewidth - 8\tabcolsep) * \real{0.2000}}
  >{\raggedright\arraybackslash}p{(\linewidth - 8\tabcolsep) * \real{0.2000}}@{}}
\toprule\noalign{}
\begin{minipage}[b]{\linewidth}\raggedright
Season
\end{minipage} & \begin{minipage}[b]{\linewidth}\raggedright
Episodes
\end{minipage} & \begin{minipage}[b]{\linewidth}\raggedright
False-suspect episodes (Y)
\end{minipage} & \begin{minipage}[b]{\linewidth}\raggedright
Y-only rate
\end{minipage} & \begin{minipage}[b]{\linewidth}\raggedright
Persons per episode
\end{minipage} \\
\midrule\noalign{}
\endhead
\bottomrule\noalign{}
\endlastfoot
1 & 21 & 10 & 47.6\% & 0.81 \\
2 & 19 & 14 & 73.7\% & 1.53 \\
3 & 23 & 18 & 78.3\% & 1.09 \\
4 & 24 & 14 & 58.3\% & 1.33 \\
5 & 23 & 18 & 78.3\% & 1.48 \\
6 & 23 & 16 & 69.6\% & 1.17 \\
7 & 22 & 11 & 50.0\% & 0.82 \\
8 & 21 & 17 & 81.0\% & 1.05 \\
9 & 18 & 15 & 83.3\% & 1.89 \\
10 & 22 & 18 & 81.8\% & 1.23 \\
11 & 23 & 18 & 78.3\% & 1.26 \\
12 & 24 & 17 & 70.8\% & 0.92 \\
13 & 23 & 16 & 69.6\% & 1.26 \\
14 & 24 & 13 & 54.2\% & 0.83 \\
15 & 24 & 13 & 54.2\% & 0.62 \\
16 & 23 & 12 & 52.2\% & 0.61 \\
17 & 23 & 11 & 47.8\% & 0.70 \\
18 & 21 & 11 & 52.4\% & 0.81 \\
19 & 24 & 13 & 54.2\% & 0.62 \\
20 & 24 & 15 & 62.5\% & 0.75 \\
21 & 20 & 6 & 30.0\% & 0.30 \\
22 & 16 & 6 & 37.5\% & 0.50 \\
23 & 22 & 9 & 40.9\% & 0.41 \\
24 & 22 & 8 & 36.4\% & 0.59 \\
25 & 13 & 5 & 38.5\% & 0.62 \\
26 & 22 & 12 & 54.5\% & 0.64 \\
27 & 3 & 1 & 33.3\% & 1.00 \\
\end{longtable}
}

\emph{Note.} Episode denominators exclude the nine corrupted-source
episodes (in Seasons 1, 2, 4, 5, 8, 9, and 11). Season 27 was three
episodes old at collection.

\textbf{Table 2. False-suspect episode rate and falsely accused persons
per episode, by era.}

{\def\LTcaptype{none} % do not increment counter
\begin{longtable}[]{@{}
  >{\raggedright\arraybackslash}p{(\linewidth - 12\tabcolsep) * \real{0.1429}}
  >{\raggedright\arraybackslash}p{(\linewidth - 12\tabcolsep) * \real{0.1429}}
  >{\raggedright\arraybackslash}p{(\linewidth - 12\tabcolsep) * \real{0.1429}}
  >{\raggedright\arraybackslash}p{(\linewidth - 12\tabcolsep) * \real{0.1429}}
  >{\raggedright\arraybackslash}p{(\linewidth - 12\tabcolsep) * \real{0.1429}}
  >{\raggedright\arraybackslash}p{(\linewidth - 12\tabcolsep) * \real{0.1429}}
  >{\raggedright\arraybackslash}p{(\linewidth - 12\tabcolsep) * \real{0.1429}}@{}}
\toprule\noalign{}
\begin{minipage}[b]{\linewidth}\raggedright
Era
\end{minipage} & \begin{minipage}[b]{\linewidth}\raggedright
Episodes
\end{minipage} & \begin{minipage}[b]{\linewidth}\raggedright
False-suspect episodes (Y)
\end{minipage} & \begin{minipage}[b]{\linewidth}\raggedright
Y-only rate
\end{minipage} & \begin{minipage}[b]{\linewidth}\raggedright
Y+Maybe rate
\end{minipage} & \begin{minipage}[b]{\linewidth}\raggedright
Persons harmed
\end{minipage} & \begin{minipage}[b]{\linewidth}\raggedright
Persons per episode
\end{minipage} \\
\midrule\noalign{}
\endhead
\bottomrule\noalign{}
\endlastfoot
Pre-\#MeToo, S1-18 (1999-spring 2017) & 401 & 262 & 65.3\% & 68.3\% &
427 & 1.06 \\
\#MeToo era, S19-21 (2017-2020) & 68 & 34 & 50.0\% & 52.9\% & 39 &
0.57 \\
Post-2020, S22-27 & 98 & 41 & 41.8\% & 42.9\% & 55 & 0.56 \\
\end{longtable}
}

\emph{Note.} Era x false suspect (Y vs.~not Y): chi-square(2, N = 567) =
20.89, exact p = .000029. Under the Y+Maybe treatment: chi-square =
24.45, exact p = .000005. Fifteen episodes carry the Maybe flag,
including four with missing transcripts.

\textbf{Table 3. Conditional consequence severity and police apology, by
era (person level, n = 521).}

{\def\LTcaptype{none} % do not increment counter
\begin{longtable}[]{@{}
  >{\raggedright\arraybackslash}p{(\linewidth - 10\tabcolsep) * \real{0.1667}}
  >{\raggedright\arraybackslash}p{(\linewidth - 10\tabcolsep) * \real{0.1667}}
  >{\raggedright\arraybackslash}p{(\linewidth - 10\tabcolsep) * \real{0.1667}}
  >{\raggedright\arraybackslash}p{(\linewidth - 10\tabcolsep) * \real{0.1667}}
  >{\raggedright\arraybackslash}p{(\linewidth - 10\tabcolsep) * \real{0.1667}}
  >{\raggedright\arraybackslash}p{(\linewidth - 10\tabcolsep) * \real{0.1667}}@{}}
\toprule\noalign{}
\begin{minipage}[b]{\linewidth}\raggedright
Era
\end{minipage} & \begin{minipage}[b]{\linewidth}\raggedright
n
\end{minipage} & \begin{minipage}[b]{\linewidth}\raggedright
Mean severity (median)
\end{minipage} & \begin{minipage}[b]{\linewidth}\raggedright
Severity 3-4
\end{minipage} & \begin{minipage}[b]{\linewidth}\raggedright
Severity 4
\end{minipage} & \begin{minipage}[b]{\linewidth}\raggedright
Any apology
\end{minipage} \\
\midrule\noalign{}
\endhead
\bottomrule\noalign{}
\endlastfoot
Pre-\#MeToo, S1-18 & 427 & 2.67 (3) & 60.0\% (256/427) & 20.1\% & 7.5\%
(32/427) \\
\#MeToo era, S19-21 & 39 & 3.05 (3) & 82.1\% (32/39) & 30.8\% & 5.1\%
(2/39) \\
Post-2020, S22-27 & 55 & 2.84 (3) & 67.3\% (37/55) & 21.8\% & 9.1\%
(5/55) \\
\end{longtable}
}

\emph{Note.} Era x severity (3-4 vs.~1-2): chi-square(2, n = 521) =
8.06, p = .018, uncorrected; the contrast does not survive a Bonferroni
correction across the four planned tests and is treated as suggestive.
Severity scale: 1 private harm, 2 public exposure without formal
sanction, 3 material sanction, 4 life-altering or death. Means treat the
scale as interval for summary only.

\textbf{Table 4. Any police threat or coercion toward innocent persons,
by era and by three-season block.}

Panel A. By era (pooled test nonsignificant: chi-square(2, n = 521) =
5.30, p = .071).

{\def\LTcaptype{none} % do not increment counter
\begin{longtable}[]{@{}llll@{}}
\toprule\noalign{}
Era & n & Any threat/coercion & Rate \\
\midrule\noalign{}
\endhead
\bottomrule\noalign{}
\endlastfoot
Pre-\#MeToo, S1-18 & 427 & 230 & 53.9\% \\
\#MeToo era, S19-21 & 39 & 14 & 35.9\% \\
Post-2020, S22-27 & 55 & 32 & 58.2\% \\
\end{longtable}
}

Panel B. By three-season block.

{\def\LTcaptype{none} % do not increment counter
\begin{longtable}[]{@{}lll@{}}
\toprule\noalign{}
Seasons & n & Rate \\
\midrule\noalign{}
\endhead
\bottomrule\noalign{}
\endlastfoot
S1-3 & 71 & 57.7\% \\
S4-6 & 93 & 59.1\% \\
S7-9 & 74 & 59.5\% \\
S10-12 & 78 & 60.3\% \\
S13-15 & 64 & 46.9\% \\
S16-18 & 47 & 27.7\% \\
S19-21 & 39 & 35.9\% \\
S22-24 & 30 & 56.7\% \\
S25-27 & 25 & 60.0\% \\
\end{longtable}
}

\emph{Note.} The dip begins at Season 15 (26.7\%); Season 16 is the
single lowest season (14.3\%). The five lowest single seasons are S16
(14.3\%), S24 (23.1\%), S15 (26.7\%), S19 (26.7\%), and S18 (29.4\%).
Rebound-era single seasons (S22-S26) range from 23.1\% to 87.5\% on
cells of 8 to 14 persons; Season 27 stands at three persons of three.
Season cells after S14 hold 3 to 18 persons; blocks are the unit of
inference.

\textbf{Table 5. Accusation origin by era (column \% within era, n =
521).}

{\def\LTcaptype{none} % do not increment counter
\begin{longtable}[]{@{}llll@{}}
\toprule\noalign{}
Origin & Pre-\#MeToo, S1-18 & \#MeToo era, S19-21 & Post-2020, S22-27 \\
\midrule\noalign{}
\endhead
\bottomrule\noalign{}
\endlastfoot
squad\_inference & 242 (56.7\%) & 16 (41.0\%) & 26 (47.3\%) \\
victim\_ID & 66 (15.5\%) & 12 (30.8\%) & 16 (29.1\%) \\
fabrication & 56 (13.1\%) & 7 (17.9\%) & 3 (5.5\%) \\
witness\_misID & 33 (7.7\%) & 2 (5.1\%) & 4 (7.3\%) \\
tech\_db\_error & 22 (5.2\%) & 0 (0.0\%) & 3 (5.5\%) \\
coerced\_interview & 7 (1.6\%) & 2 (5.1\%) & 3 (5.5\%) \\
unknown & 1 (0.2\%) & 0 (0.0\%) & 0 (0.0\%) \\
\textbf{n} & \textbf{427} & \textbf{39} & \textbf{55} \\
\end{longtable}
}

\emph{Note.} Era n for S19-21 and S22-27 are 39 and 55; swings under
roughly ten points should be treated as noise.

\textbf{Table 6. Media exposure of falsely accused persons, by era.}

{\def\LTcaptype{none} % do not increment counter
\begin{longtable}[]{@{}llll@{}}
\toprule\noalign{}
Era & n & Media channel & Media + online \\
\midrule\noalign{}
\endhead
\bottomrule\noalign{}
\endlastfoot
Pre-\#MeToo, S1-18 & 427 & 65 (15.2\%) & 66 (15.5\%) \\
\#MeToo era, S19-21 & 39 & 13 (33.3\%) & 13 (33.3\%) \\
Post-2020, S22-27 & 55 & 8 (14.5\%) & 9 (16.4\%) \\
\end{longtable}
}

\emph{Note.} Full exposure-channel distributions by era, including
workplace, family, school, legal, and church channels, are available
from the author; only the media channel moves appreciably across eras.

\end{document}
